\section{Introduction}

A language model's alignment strategy --- whether it was fine-tuned with Direct Preference
Optimization (DPO) or Supervised Fine-Tuning (SFT) --- can be inferred from four standard
benchmark scores alone, without access to training data, preference datasets, or model
weights, at 83.3\% accuracy. The dimension that drives this inference is not fairness,
as alignment researchers might expect, but truthfulness.

This finding has immediate practical stakes. As language models are deployed at scale,
practitioners and auditors increasingly need to verify how models were trained --- yet
alignment provenance is frequently undocumented, especially for community-released models.
A benchmark-based fingerprinting method would enable alignment auditing from public
leaderboard data, requiring no whitebox access to training pipelines. But such a method
is only trustworthy if we understand which benchmark dimensions carry alignment-strategy
signal and which do not.

The dominant narrative surrounding DPO alignment, as articulated in the original DPO
paper \cite{Rafailov2023Direct} and subsequent commentary, holds that preference-based
training rewards annotator-preferred responses. Since human annotators penalize biased
and harmful content, DPO is expected to produce a measurable fairness advantage over
SFT. Accordingly, one would predict that fairness benchmarks --- bias-eliciting question
answering (BBQ), gender pronoun resolution (WinoGender) --- should be the primary
dimensions separating DPO and SFT model profiles. This intuition has shaped how
researchers think about the behavioral consequences of DPO alignment.

However, this narrative has a critical empirical gap: it has never been tested against
matched DPO/SFT controls using a multi-dimensional trustworthiness benchmark suite.
Prior work evaluates DPO on capability benchmarks (MT-Bench, AlpacaEval) \cite{Rafailov2023Direct},
or evaluates trustworthiness across model families without isolating alignment
strategy \cite{Wang2023DecodingTrust,Liang2022Holistic}. No study has systematically compared
DPO and SFT models on unified trustworthiness benchmarks while controlling for base
model and data, or applied a classification framing to test whether alignment strategy
is detectable from benchmark profiles.

The gap this creates is concrete: without empirical characterization of which benchmark
dimensions carry alignment-strategy signal, any alignment auditing tool built on assumed
fairness-DPO correspondence risks instrumentalizing the wrong dimension --- or failing
to detect alignment differences at all.

Our key insight is that alignment strategy can be treated as a binary classification
problem over public benchmark score vectors. A model's scores on TruthfulQA MC2, BBQ,
WinoGrande, and WinoGender form a point in 4D space; DPO and SFT models cluster into
distinguishable regions in that space. A nearest-neighbor classifier can recover a new
model's likely alignment strategy at statistically significant accuracy --- and the
dimension driving the separation reveals which aspect of behavior preference training
actually modifies most systematically.

We make the following contributions:

\textbf{C1 (Empirical, Positive): Alignment fingerprinting is feasible from public benchmarks.}
Across 12 7B-parameter models (6 DPO, 6 SFT), a k-NN (k=1) leave-one-out classifier
achieves 83.3\% accuracy (10/12 correct; permutation p=0.031). Alignment strategy is
recoverable from a 4D benchmark profile with no whitebox model access required.

\textbf{C2 (Empirical, Surprising): Truthfulness, not fairness, drives the alignment fingerprint.}
Fisher's discriminant criterion reveals that TruthfulQA MC2 carries 9.5$\times$ more class-separability
signal than the next benchmark (Fisher=0.8122 vs WinoGender=0.0856).
DPO models score +4.6pp higher on truthfulness on average across matched pairs
(per-pair mean delta) --- a direction that contradicts the dominant theoretical prediction. BBQ fairness shows no systematic DPO advantage
(k=3/6, p=0.66), refuting the bias-avoidance mechanism as the primary driver.

\textbf{C3 (Methodological): Alignment detection as a classification task.}
We introduce the framing of alignment-strategy fingerprinting as a binary classification
problem over benchmark score vectors, providing a practical tool for model auditing and
provenance verification. Misclassification cases are mechanistically informative: an
RLHF model (Llama-2-chat) clusters with DPO models, suggesting preference-based methods
share a benchmark signature distinct from pure SFT; a DPO model trained on the same
base as its SFT counterpart (zephyr pair) is misclassified, showing that base architecture
can dominate alignment signal for near-identical pairs.

\textbf{C4 (Implication): The DPO fairness assumption requires empirical re-examination.}
Our results --- to our knowledge, the first paired DPO/SFT comparison on a unified fairness-inclusive
trustworthiness suite --- provide no support for a systematic DPO fairness advantage.
This challenges a widely-held intuition about what preference-based training optimizes
and motivates revisiting the role of annotation criteria in shaping DPO's behavioral
profile.

We proceed as follows. Section~\ref{sec:related} reviews related work on alignment evaluation,
trustworthiness benchmarking, and model fingerprinting. Section~\ref{sec:method} describes our
experimental methodology, including model selection, benchmark suite, and classification
approach. Section~\ref{sec:setup} presents experimental design details. Section~\ref{sec:results} reports results
for both fingerprint existence (h-e1) and mechanism analysis (h-m1). Section~\ref{sec:discussion} discusses
interpretations, limitations, and scope conditions. Section~\ref{sec:conclusion} concludes.
